\documentstyle[%


righttag%


,11pt%


,amssymb%


,russian%               remove in Eng.


,izvru%                 Set izven% in Eng.


]{amsart}





%





\newcommand{\id}{\operatorname{id}}


\newcommand{\tts}[1]{}





\theoremstyle{definition}


\theorembodyfont{\rm}


\newtheorem{notenonum}{\hskip\parindent Замечание}


\renewcommand{\thenotenonum}{}





%\hoffset=-15mm%                  For Eng.


\setcounter{page}{71}


\god{2005}


\nomer{6~(517)} %                        Remove in Eng.


%\nomer{Vol.\,49, No.\,6}%               Set in Eng.


%\str{\pageref{begin}--\pageref{end}}%  Set in Eng.


\udk{515.124}





\author{\it Е.Н.\,Сосов}


% NB:   ^^ remove in Eng.


\title{Касательное пространство по Буземану}





\date{}


%\pagestyle{myheadings}%         Set in Eng.


\pagestyle{plain} %              Remove in Eng.


\begin{document}


%\thispagestyle{plain}%          Set in Eng.


\maketitle


\label{begin}








Исследуются метрические свойства


касательного пространства для метрического пространства более


общего, чем дифференцируемое $G$-пространство Буземана. Установлено,


что метрика на касательном пространстве в произвольной точке пространства


неположительной кривизны по Буземану (дифференцируемого по Буземану


метрического пространства) внутренняя. Доказано, что касательное


пространство в произвольной точке локально полного дифференцируемого


по Буземану метрического пространства является полным, а


также, что касательное пространство в произвольной точке локально


компактного пространства неположительной  кривизны по Буземану


является конечно-компактным геодезическим.








\section{Необходимые определения и теоремы}





Рассмотрим метрическое пространство $(X,\rho)$ с выделенным


семейством сегментов $S$ (сегментом $[x, y]$ с концами $x, y \in


X$ называется непрерывная кривая, соединяющая точки $x$, $y$,


длина которой равна расстоянию $\rho(x,y)$ между этими точками


([1], с.\,42)), удовлетворяющее следующим условиям.


\begin{itemize}


\item[A.] Каждый подсегмент произвольного сегмента из $S$ принадлежит $S$.


\item[B.] Для каждого $p \in X$ найдется такое положительное


вещественное число $r(p)$, что для любых $x$, $y$ из открытого


шара $B(p, r(p))$ существует единственный сегмент


$[x, y] \in S$.


\end{itemize}





В дальнейшем будем использовать следующие обозначения и определения.


${\Bbb R}_{+}$ --- множество неотрицательных


вещественных чисел; $xy=\rho(x,y)$.


$B[x,r]$ ($B(x,r)$, $S(x,r)$) --- замкнутый шар


(открытый шар, сфера) с центром в точке $x \in X$, радиуса $r >0$.


Для $\lambda \in [0,1]$,


$x, y \in X$ \ $\omega_{\lambda}[x,y]$ --- точка сегмента $[x,y] \in S$


такая, что $x\omega_{\lambda}[x,y]=\lambda xy$.





Пространство $(X,\rho)$, удовлетворяющее условиям A, B, назовем


дифференцируемым в точке $p \in X$ по Буземану, если для каждого


$\varepsilon>0$ найдется такое


$\delta \in (0, r(p))$,


что $|\omega_{\lambda}[p, x]\omega_{\lambda}[p,y]-\lambda xy|\leq


\varepsilon \lambda xy$ 


для любых $\lambda \in [0, 1]$, $x, y \in B(p,\delta)$


([1], с.\,293; [2], с.\,21).





Пространство $(X,\rho)$, удовлетворяющее условиям A, B, назовем


пространством неположительной кривизны по Буземану, если


$2\omega_{1/2}[z,x]\omega_{1/2}[z,y] \leq xy$ для любых


$x, y, z \in B(p, r(p))$ ([1], с.\,304; [3], c.\,63).





Метрическое пространство $X$ называется пространством с внутренней метрикой,


если для любых $x,y\in X$, $\varepsilon>0$ найдется конечная 


последовательность точек


$z_{0}=x$, $z_{1},\dotsb, z_{k}=y$ такая, что $z_{i}z_{i+1}<\varepsilon$


($0 \leq i \leq k-1$) и $z_{0}z_{1}+\dotsb+z_{k-1}z_{k}<xy+\varepsilon$ [4].





Пространство $(X,\rho)$ называется геодезическим,


если любые две его точки можно соединить сегментом [5] (в


первоначальном варианте в этом определении пространство $X$


считали полным метрическим [6]).





Пространство $(X,\rho)$ называется конечно-компактным метрическим


пространством, если в нем каждое ограниченное замкнутое множество


компактно ([1], с.\,18).





Отображение $f$ из метрического пространства $(X,\rho)$ в


метрическое пространство $(Y,d)$ называется билипшицевым


отображением, если найдется такая константа


$\lambda \geq 1$,


что $\rho(x,y)/\lambda \leq d(f(x),f(y)) \leq \lambda\rho(x,y)$


для любых $x,y \in X$.


Метрические пространства $X$, $Y$ называются билипшицево


эквивалентными, если существует сюръективное билипшицево


отображение $f : X \rightarrow Y$ ([7], с.\,269).





Сформулируем полученные результаты.








\begin{lem}


Пусть пространство $(X,\rho)$ удовлетворяет


условиям {\em A, B}. Тогда функция


$$


\overline{m}_{p}: B(p, r(p))\times B(p, r(p))


\longrightarrow {\Bbb R}_{+}, \quad \overline{m}_{p}(x, y)=


\underset{\lambda \to 0+}{\varlimsup}{}{\frac{\omega_{\lambda}[p,


x]\omega_{\lambda}[p, y]}{\lambda}}


$$


есть псевдометрика,


обладающая следующими свойствами\/${:}$


$$


\overline{m}_{p}(p, x)=px,\quad


\overline{m}_{p}(\omega_{\lambda}[p, x],\omega_{\lambda}[p, y])=


\lambda \overline{m}_{p}( x, y)


\ \ \text{для любых $\lambda \in [0, 1]$},


\ \ x, y \in B(p, r(p)).


$$


\end{lem}





Пусть пространство $(X,\rho)$ удовлетворяет условиям A, B. Из


леммы 1 следует, что функция


$m_{p}: (B(p, r(p)) \times {\Bbb R}_{+}) \times (B(p, r(p)) \times


{\Bbb R}_{+}) \longrightarrow {\Bbb R}_{+}$,


$m_{p}((x;\lambda),(y;\mu))=


\tau\overline{m}_{p}(\omega_{\lambda/\tau}[p,x],\omega_{\mu/\tau}[p,y])$,


где $\tau \in [\max[\lambda,\mu],+\infty)$, --- корректно определенная


псевдометрика. Введем на множестве $B(p, r(p)) \times {\Bbb


R}_{+}$ отношение эквивалентности


$(x;\lambda)\sim (y;\mu)$, если


$m_{p}((x;\lambda),(y;\mu))=0$. Тогда псевдометрика индуцирует


на фактор-пространстве  $X_{p}=(B(p, r(p)) \times {\Bbb


R}_{+})/\sim$ метрику, которую обозначим тем же символом $m_{p}$.








\begin{thm}


Пусть $(X,\rho)$ --- пространство неположительной кривизны по Буземану. Тогда


\begin{itemize}


\item[(i)]


$\overline{m}_{p}(x, y)=\lim \limits_{\nu \to 0+}


{\frac{\omega_{\nu}[p, x]\omega_{\nu}[p,y]}{\nu}}$


для всех $x, y\in B(p,r(p));$


\item[(ii)] $\overline{m}_{p}(x, y) \leq xy$ для всех $x, y \in B(p,r(p));$


\item[(iii)] $m_{p}$ есть внутренняя метрика на $X_{p};$


\item[(iv)] если, кроме того, пространство $X$ локально компактно, то для


любого $p \in X$ \ $(X_{p},m_{p})$ --- конечно-компактное геодезическое


пространство.


\end{itemize}


\end{thm}








\begin{thm}


Пусть $(X,\rho)$ --- пространство, дифференцируемое в


точке $p \in X$ по Буземану. Тогда


\begin{itemize}


\item[(i)]


$\overline{m}_{p}(x, y)=\lim \limits_{\nu \to 0+}


{\frac{\omega_{\nu}[p, x]\omega_{\nu}[p,y]}{\nu}}$


для всех $x, y\in B(p,r(p));$


\item[(ii)] $\overline{m}_{p}$ есть метрика на открытом шаре $B(p, r(p))$,


и для каждого $\varepsilon>0$ найдется такое $\delta \in (0,


r(p))$, что $|\overline{m}_{p}(x, y) - xy| \leq \varepsilon xy$


для всех $x, y \in B(p,\delta);$


\item[(iii)] $m_{p}$ есть внутренняя метрика на $X_{p};$


\item[(iv)] если, кроме того, $X$ --- локально полное пространство, то


$(X_{p},m_{p})$ --- полное пространство\/$;$


\item[(v)] пусть $\delta$ --- то же, что и в п.\,{\em (ii)} при $0<\varepsilon<1$,


и пусть найдется такое число $t(p) \in (0,\delta)$, что для


каждой точки $x \in B(p,t(p))\setminus\{p\}$ существует точка


$\wh x \in S(p,t(p))$


такая, что $x \in [p,\wh x]$, где $[p,\wh x] \in S$, тогда для каждого


$R>0$ пространства $(B[p,t(p)],\rho)$, $(B[[p;1],R],m_p)$ билипшицево


эквивалентны.


\end{itemize}


\end{thm}








\begin{notenonum}


Нетрудно проверить, что если $X$ --- дифференцируемое в


точке $p$ \ $G$-про\-ст\-ран\-ст\-во Буземана, то пространство $(X_{p},m_{p})$ 


изометрично касательному пространству в точке $p$ по Буземану (при наличии в


пространстве $(X_{p},m_{p})$ линейной структуры эта изометрия


является линейной) ([2], c.\,22). Поэтому метрическое пространство


$(X_{p},m_{p})$ естественно называть касательным пространством по


Буземану в точке $p \in X$.


\end{notenonum}








\section{Доказательства полученных результатов}





{\bf Доказательство леммы 1} сводится к непосредственной проверке.


\medskip





\begin{pf*}{Доказательство теоремы 1}


(i)--(ii) Пусть $(X,\rho)$ удовлетворяет условиям A, B и\\


$2\omega_{1/2}[z, x]\omega_{1/2}[z, y]


\leq xy$ для любых $x, y, z \in B(p, r(p))$.


Из последнего вытекает


\begin{equation}


\omega_{\lambda}[z,x]\omega_{\lambda}[z,y] \leq \lambda xy


\end{equation}


для любых $0 \leq \lambda \leq 1$, $x, y, z \in


B(p,r(p))$. Следовательно, $\overline{m}_{p}( x, y) \leq xy$ для


любых $x, y \in B(p, r(p))$. Кроме того, из доказательств теорем


36.4--36.6 в ([1], с.\,304--307) следует, что для любых $x, y, z


\in B(p,r(p))$ непрерывная функция $f : [0,1] \rightarrow R$,


$f[\lambda ]=\omega_{\lambda}[z,x]\omega_{\lambda}[z,y]$


выпукла. Таким образом, существует предел


$\lim\limits_{\nu \to 0+}{\frac{\omega_{\nu}[p, x]\omega_{\nu}[p,y]}{\nu}}$


для $x, y \in B(p,r(p))$.





(iii) Выберем произвольно $[x_1;\lambda], [y;\mu] \in X_p$ и для


определенности считаем, что $\lambda \leq \mu$. Тогда


$[x_1;\lambda]=[\omega_{\lambda/\mu}[p,x_1];\mu]$. Обозначим $x


=\omega_{\lambda/\mu}[p,x_1]$. В силу леммы 1 из [8] достаточно


доказать, что для каждого $\varepsilon >0$ найдется такая точка


$[z;\tau ] \in X_p$, что $2\max[m_p([x;\mu],[z;\tau


]),m_p([y;\mu],[z;\tau ])]<m_p([x;\mu],[y;\mu])+\varepsilon$.


Очевидно, найдется такое $\nu_0 \in (0,1]$, что $z_\nu=


\omega_{1/2}[\omega_\nu[p,x],\omega_\nu[p,y]] \in B(p,r(p))$ для


$\nu \in (0,\nu_0]$. 


Из неравенства (1) следует 


\begin{equation}


\max[\overline{m}_p(\omega_\nu[p,x],z_\nu),\overline{m}_p


(\omega_\nu[p,y],z_\nu)]\leq 


\max[\omega_\nu[p,x]z_\nu,\omega_\nu[p,y]z_\nu]=


\omega_\nu[p,x]\omega_\nu[p,y]/2.


\end{equation}


Отсюда и из леммы 1 получим


$\nu\overline{m}_p(x,y)=


\overline{m}_p(\omega_\nu[p,x],\omega_\nu[p,y]) \leq


\overline{m}_p(\omega_\nu[p,x],z_\nu)+


\overline{m}_p(\omega_\nu[p,y],z_\nu) \leq


\omega_\nu[p,x]\omega_\nu[p,y]$.


Из этих неравенств и утверждения (i) вытекает


$$


\lim\limits_{\nu \to 0+}{\frac{\overline{m}_p(\omega_\nu[p,x],z_\nu)+


\overline{m}_p(\omega_\nu[p,y],z_\nu)}{\nu}}=\overline{m}_p(x,y).


$$


Если учтем (2), то получим 


\begin{equation}


\lim\limits_{\nu \to 0+}{\frac{\overline{m}_p(\omega_\nu[p,x],z_\nu)}{\nu}}


=\lim\limits_{\nu \to 0+}{\frac{\overline{m}_p(\omega_\nu[p,y],z_\nu)}{\nu}}=


\overline{m}_p(x,y)/2.


\end{equation}


Следовательно, для каждого


$\varepsilon>0$ найдется такое $\nu \in (0,\min[\delta,1])$, что


$$


\frac{\overline{m}_p(\omega_\nu[p,x],z_\nu)}{\nu} <\overline{m}_p(x,y)/2


+\varepsilon, \quad \frac{\overline{m}_p(\omega_\nu[p,y],z_\nu)}{\nu} <


\overline{m}_p(x,y)/2+\varepsilon.


$$


Тогда для каждого


$\varepsilon>0$ найдем такое $\nu \in (0,\min[\delta,1])$,


что 


\begin{gather*}


2\max[m_p([x;\mu],[z_\nu;\mu/\nu]),m_p([y;\mu],[z_\nu;\mu/\nu])]= \\


%


=2\mu/\nu\max[\overline{m}_p(\omega_\nu[p,x],z_\nu),


\overline{m}_p(\omega_\nu[p,y],z_\nu)]


< \mu\overline{m}_p(x,y)+2\mu\varepsilon=m_p([x;\mu],[y;\mu])+


2\mu\varepsilon.


\end{gather*}





(iv) Пусть $X$ --- локально-компактное


пространство неположительной кривизны по Буземану и


$\{[x_n;\lambda_n]\}$ --- произвольная ограниченная


последовательность пространства $(X_{p},m_{p})$. Заметим, что


$m_p ([x_n;\lambda_n],[p;1])=


\tau\overline {m}_{p}(\omega_{\lambda_n/\tau}[p,x_n],p)=\lambda_n px_n$


для $\tau \geq \max [1,\lambda_n]$


и последовательность $\{\lambda_n\}$ ограничена. Тогда найдутся


такие вещественные числа $r \in (0,r(p)]$, $\nu>0$, что шар


$B[p,r]$ компактен и $\omega_{\lambda_n/\nu}[p,x_n] \in B[p,r]$


($n=1,2,\dotsc$). Выделим подпоследовательность


$\{\omega_{\lambda_k/\nu}[p,x_k]\} \subset


\{\omega_{\lambda_n/\nu}[p,x_n]\}$, сходящуюся в исходной метрике


к точке $y \in B[p,r]$. Отсюда, из леммы 1 и неравенства (1)


следует, что $m_p ([x_k;\lambda_k],[y;\nu])=


\nu\overline{m}_p(\omega_{\lambda_k/\nu}[p,x_k],y) \leq


\nu\omega_{\lambda_k/\nu}[p,x_k]y \rightarrow 0$ ($k \rightarrow \infty$).


Таким образом, последовательность $\{[x_n;\lambda_n]\}$ обладает


сходящейся подпоследовательностью, и пространство


$(X_{p},m_{p})$ конечно-компактно. Из следствия обобщенной


альтернативы Хопфа--Ринова ([9], с.\,297) и конечной компактности


пространства $(X_{p},m_{p})$ следует, что оно 


геодезическое.


\end{pf*}








\begin{pf*}{Доказательство теоремы 2} 


(i)--(ii) Для каждого $\varepsilon>0$


найдется такое


$\delta(\varepsilon) \in (0, r(p))$, что для любых $x, y \in


B(p,\delta(\varepsilon))$, $0<\lambda<\mu<1$ верно неравенство


$$


\bigg|\omega_{\lambda}[p,x]\omega_{\lambda}[p,y] -


\frac{\lambda}{\mu} \omega_{\mu}[p,x]\omega_{\mu}[p,y] \bigg|\leq \varepsilon


\frac{\lambda}{\mu}\omega_{\mu}[p,x]\omega_{\mu}[p,y].


$$


Тогда для любых $x, y \in B(p,\delta[1/2])$, $0<\mu<\mu_0 <


\frac{\delta[1/2]}{r(p)}$ верны неравенства


$$


\frac{\omega_{\mu_0}[p,x]\omega_{\mu_0}[p,y]}{2\mu_0}\leq


\frac{\omega_{\mu}[p,x]\omega_{\mu}[p,y]}{\mu} \leq


\frac{3\omega_{\mu_0}[p,x]\omega_{\mu_0}[p,y]}{2\mu_0},


$$


и для любых $0<\varepsilon<1/2$,


$x, y \in B(p,\delta(\varepsilon))$,


$\mu<\min [\frac{\delta(\varepsilon)}{r(p)},\mu_0]$ выполняются


неравенства


$$


\bigg|\frac{\omega_{\lambda}[p,x]\omega_{\lambda}[p,y]}{\lambda } -


\frac{\omega_{\mu}[p,x]\omega_{\mu}[p,y]}{\mu} \bigg|\leq \varepsilon


\frac{\omega_{\mu}[p,x]\omega_{\mu}[p,y]}{\mu} \leq


\frac{3\varepsilon \omega_{\mu_0}[p,x]\omega_{\mu_0}[p,y]}{2\mu_0}.


$$


Отсюда следует существование предела


$$


\overline{m}_{p}( x, y)=


\lim\limits_{\nu \to 0+}{\frac{\omega_{\nu}[p, x]\omega_{\nu}[p,


y]}{\nu}}


$$


и то, что $\overline{m}_{p}( x, y) \neq 0$ при $x \neq


y$. Если обе части неравенства из определения дифференцируемости в


точке $p \in X$ по Буземану поделить на $\lambda>0$ и перейти к


пределу при $\lambda \to 0+$, то получим неравенство


$|\overline{m}_{p}( x, y) - xy| \leq \varepsilon xy$ для


$x, y \in B(p,\delta)$.





(iii) Выберем произвольно $[x_1;\lambda], [y;\mu] \in X_p$ и для


определенности положим $\lambda \leq \mu$. Тогда


$[x_1;\lambda]=[\omega_{\lambda/\mu}[p,x_1];\mu]$. Обозначим $x


=\omega_{\lambda/\mu}[p,x_1]$. В силу леммы 1 из [8] достаточно


доказать, что для каждого $\varepsilon >0$ найдется такая точка


$[z;\tau ] \in X_p$, что $2\max[m_p([x;\mu],[z;\tau


]),m_p([y;\mu],[z;\tau ])]<m_p([x;\mu],[y;\mu])+\varepsilon$.


Очевидно, найдется такое $\nu_0 \in (0,1]$, что $z_\nu=


\omega_{1/2}[\omega_\nu[p,x],\omega_\nu[p,y]] \in B(p,r(p))$ для


$\nu \in (0,\nu_0]$. В силу утверждения (ii) для каждого


$\varepsilon>0$ найдутся такие $\delta>0$ и


$\nu_0\in (0,1]$, что $\omega_\nu[p,x], \omega_\nu[p,y],


\omega_{1/2}[\omega_\nu[p,x],\omega_\nu[p,y]] \in B(p,\delta)$


и


$\overline{m}_p(\omega_\nu[p,x],\omega_\nu[p,y]) \leq (1+


\varepsilon)\omega_\nu[p,x]\omega_\nu[p,y]$ для любого $\nu \in (0,\nu_0]$.


Тогда для любых $\varepsilon>0$, $\nu \in (0,\nu_0]$


\begin{equation}


\max[\overline{m}_p(\omega_\nu[p,x],z_\nu),


\overline{m}_p(\omega_\nu[p,y],z_\nu])\leq 


(1+\varepsilon)\max[\omega_\nu[p,x]z_\nu,\omega_\nu[p,y]z_\nu]=(1+


\varepsilon)\omega_\nu[p,x]\omega_\nu[p,y]/2.


\end{equation}


Отсюда и из леммы 1 получим $\nu\overline{m}_p(x,y)=


\overline{m}_p(\omega_\nu[p,x],\omega_\nu[p,y]) \leq


\overline{m}_p(\omega_\nu[p,x],z_\nu)+


\overline{m}_p(\omega_\nu[p,y],z_\nu) \leq (1+


\varepsilon)\omega_\nu[p,x]\omega_\nu[p,y]$ для любых $\varepsilon


> 0$, $\nu \in (0,\nu_0]$.


Если учтем (4), то получим


\begin{gather*}


\overline{m}_p(x,y) \leq


\underset{\nu \to 0+}{\varlimsup}{}{\frac{\overline{m}_p(\omega_\nu[p,x],z_\nu)+


\overline{m}_p(\omega_\nu[p,y],z_\nu)}{\nu}} \leq (1+\varepsilon)


\overline{m}_p(x,y), \\


%


\underset{\nu \to 0+}{\varlimsup}{}


{\frac{\max[\overline{m}_p(\omega_\nu[p,x],z_\nu),


\overline{m}_p(\omega_\nu[p,y],z_\nu])}


{\nu}} \leq (1+\varepsilon)\overline{m}_p(x,y)/2


\end{gather*}


для любого $\varepsilon>0$. Аналогичные неравенства верны и для нижних


пределов. В силу произвольности $\varepsilon>0$


$$


\lim\limits_{\nu \to


0+}{\frac{\overline{m}_p(\omega_\nu[p,x],z_\nu)}{\nu}}=


\lim\limits_{\nu \to 0+}{\frac{\overline{m}_p(\omega_\nu[p,y],z_\nu)}{\nu}}=


\overline{m}_p(x,y)/2.


$$


Оставшаяся часть доказательства


совпадает с частью доказательства утверждения (iii) теоремы 1


после равенств (3).





(iv) Пусть $(X,\rho)$ --- локально полное дифференцируемое в точке


$p \in X$ по Буземану пространство и $(B[p,r],\rho)$, где $0<r<


\delta$, --- полное подпространство пространства $(X,\rho)$. Из неравенств


$(1-\varepsilon)xy \leq \overline{m}_{p}( x, y) \leq (1+\varepsilon)xy$ для


$x, y \in B[p,r]$, $0<\varepsilon<1$, полученных в п.\,(ii),


следует билипшицева эквивалентность пространств


$(B[p,r],\overline{m}_{p})$, $(B[p,r],\rho)$ и полнота пространства


$(B[p,r],\overline{m}_{p})$. Пусть


$\{[x_n;\lambda_n]\}$ --- фундаментальная последовательность в


пространстве $(X_{p},m_{p})$. Тогда в силу ограниченности


фундаментальной последовательности


$\{[x_n;\lambda_n]\}$ и леммы 1 найдется такое число $\tau >


0$, что $\omega_{\lambda_n/\tau}[p,x_n] \in


(B[p,r],\overline{m}_p)$ и $\{\omega_{\lambda_n/\tau}[p,x_n]\}$


является фундаментальной последовательностью. Пусть эта


последовательность сходится к точке $z \in B[p,r]$. Тогда


$m_{p}([x_n;\lambda_n],[z;\tau])=\tau \overline


{m}_{p}(\omega_{\lambda_n/\tau}[p,x_n],z) \rightarrow 0$ $(n


\rightarrow \infty)$. Таким образом, метрическое пространство


$(X_{p},m_{p})$ полное.





(v) Докажем, что искомую билипшицеву эквивалентность определяет


следующая композиция


$h\circ f\circ \id : (B[p,t(p)],\rho) \rightarrow (B[[p;1],R],m_p)$


отображений


$\id : (B[p,t(p)],\rho) \rightarrow (B[p,t(p)],\overline {m}_p)$;


$f:(B[p,t(p)],\overline {m}_p) \rightarrow (B[[p;1],t(p)],m_p)$,


$f(x)=[x;1]$; $h : (B[[p;1],t(p)],m_p) \rightarrow


(B[[p;1],R],m_p)$, $h([x;\lambda])=[x;R\lambda/t(p)]$.


Рассмотрим каждое из этих отображений. Тождественное отображение


$\id : (B[p,t(p)],\rho) \rightarrow (B[p,t(p)],\overline {m}_p)$ определяет,


как было отмечено в начале доказательства п.\,(iv), билипшицеву


эквивалетность соответствующих пространств. Докажем, что $f$ есть


изометрия. В силу равенств


$m_p(f(x),f(y))=m_p([x;1],[y;1])=\overline {m}_p(x,y)$ для всех


$x$, $y \in B[p,t(p)]$, отображение $f$ есть изометрическое вложение.


Осталось доказать, что оно сюръективно. Выберем произвольно


$[x;\mu] \in B[[p;1],t(p)]$. Если $\mu \leq 1$, то $f(\omega_{\mu}[p,x])=


[\omega_{\mu}[p,x];1]=[x;\mu]$. Если $\mu>1$ и $\lambda=\mu px/t(p)$,


то $f(\omega_{\lambda}[p,\wh x])=[\omega_{\lambda}[p,\wh x];1]=


[x;\mu]$. Следовательно, отображение $f$ сюръективно. Отображение


$h$ есть подобие (определение см. в [1], с.\,286). Действительно, отображение


$h$, очевидно, обратимо и $m_p(h([x;\lambda]),h([y;\mu]))=


m_p([x;R\lambda/t(p)],[y;R\mu/t(p)])=


(R/t(p))m_p([x;\lambda],[y;\mu])$


для любых $[x;\lambda]$, $[y;\mu] \in B[[p;1],t(p)]$.


Следовательно, отображение


$h\circ f\circ \id$ определяет билипшицеву эквивалентность пространств


$(B[p,t(p)],\rho)$, $(B[[p;1],R],m_p)$.


\end{pf*}








\begin{thebibliography}{9}





\bibitem{1} Буземан Г. {\em Геометрия геодезических}. -- М.: Физматгиз,


1962. -- 503~с.





\bibitem{2} Busemann H. {\em Recent synthetic differential geometry}.


-- Berlin--Heidelberg--New~York: Springer-Verlag, 1970. -- 110 p.





\bibitem{3} Busemann H., Phadke B.B. {\em Spaces with distinguished


geodesics}. -- New~York--Basel: Marsel Dekker Inc., 1987. -- 159~p.





\bibitem{4} Бураго Ю.Д., Громов М.Л., Перельман Г.Д.


{\em Пространства А.Д.\,Александрова с ограниченными снизу кривизнами} //


УМН. -- 1992. -- Т.\,47. -- Вып.\,2. -- С.\,3--51.





\bibitem{5} Берестовский В.Н. {\em Пространства Буземана ограниченной


сверху кривизны по Александрову} // Алгебра и анализ. -- 2002. --


Т.\,14. -- Вып.\,5. -- С.\,3--18.





\bibitem{6} Ефремович В.А. {\em Неэквиморфность пространств Эвклида и 


Лобачевского} // УМН. -- 1949. -- Т.\,4. -- Вып.\,2. -- С.\,178--179.





\bibitem{7} Bonk M., Schramm O. {\em Embeddings of Gromov hyperbolic spaces}


// Geom. funct. anal. -- 2000. -- V.\,10. -- P.\,266--306.





\bibitem{8} Sosov E.N. {\em On Hausdorff intrinsic metric}


// Lobachevskii J. of


Math. -- 2001. -- V.\,8. -- P.\,185--189.





\bibitem{9} Кон-Фоссен С.Э. {\em Некоторые вопросы дифференциальной геометрии


в целом}. -- М.: Физматгиз, 1959. -- 304~с.





\end{thebibliography}





\vskip 2cm





\post{\it Казанский государственный}{\it Поступила}


\post{\it университет}{20.10.2003}





\label{end}


\end{document}


